% -------------------------------------------------------------------------
% ÜBERHOLT (31.08.2026). Maßgeblich ist die eingereichte main.tex. Zahlen,
% Zeilennummern und Angaben zum Overleaf-Stand in dieser Datei stammen aus
% der Zeit VOR der Crosswalk-Korrektur vom 29.08.2026 (35 statt 36 Stadt-
% teile, vollständiger Neulauf) und wurden bewusst NICHT nachgezogen: die
% Datei ist als datierter Entwurf Teil der Arbeitsdokumentation.
% -------------------------------------------------------------------------
% =========================================================================
% KAPITEL 5.1 -- DATENAUSWAHL UND BEREINIGUNG
% Fassung 3, 24.08.2026. STAND IN OVERLEAF -- eingetippt und validiert.
%
% GEAENDERT GEGENUEBER FASSUNG 2
%   - Einstieg nach dem Bauplan von 5.2: Befund aus Kapitel 4 mit Zahl ->
%     was hier passiert -> warum -> wo es hinterlegt ist -> Ergebnis mit
%     Zahl. Kein "dieser Abschnitt" mehr.
%   - Schlusssatz nennt NUR den Rahmen (35 Stadtteile, 132 Monate). Die
%     Fassung mit "719.989 Einsatzmeldungen UND 35 Stadtteile ueber 132
%     Monate" war falsch: 719.989 ist der entdoppelte Gesamtbestand ab 2003
%     ueber alle 41 Stadtteile, im Rahmen liegen 350.481. Zwei Zahlen mit
%     verschiedener Bezugsmenge nebeneinander.
%   - Absatz 3 nennt 719.989 nicht mehr -- die Zahl gehoert nach 5.2, wo sie
%     als Ergebnis der Zusammenfuehrung steht.
%   - Absatz 4 nennt die 41 nicht mehr, sie steht jetzt in Absatz 1.
%
% GEAENDERT GEGENUEBER FASSUNG 1
%   - Absatz 2 zaehlt nicht mehr auf, was fehlt, sondern sagt positiv, was
%     aufgenommen wird und warum. Drei Beispiele statt zehn Nennungen.
%   - Abbildung A0 entfaellt.
%
% WAS AUS FASSUNG 1 VERSCHWUNDEN IST UND WOHIN ES GEHT
%   Die beiden verworfenen Aggregationsebenen (Einzeleinsatz, Census Tract
%   und Jahr) stehen jetzt nirgends in 5.1. Sie gehoeren nach 5.3, wo die
%   Klassifikationsebene mit einer gemessenen Zahl begruendet wird
%   (350.481 Zeilen, 4.619 Profile, Obergrenze 49,9 gegen 48,2 Prozent).
%   Dort sind sie ein Befund statt einer Behauptung. Im Geruest nachgezogen.
%
% VOR DEM EINFUEGEN
%   1) DEN ALTEN FLIESSTEXT LOESCHEN, der heute zwischen
%      \section{Data Preparation} und \subsection{Datenauswahl und
%      Bereinigung} steht -- alle 13 Zeilen, samt \footcite{Bergmeir2012}.
%      Er wird durch diesen Abschnitt ersetzt.
%   2) LABELS SIND ERLEDIGT -- am 24.08.2026 direkt in Overleaf gesetzt.
%      Kapitel 4 hatte gar keine Labels, 2.6 auch nicht. Neu eingefuegt:
%        Zeile 772  \label{sec:modellbewertung}  nach
%                   \subsection{Modelloptimierung und Modellbewertung}  (2.6)
%        Zeile 894  \label{sec:datengrundlage}   nach
%                   \subsection{Datengrundlage und Merkmale}            (4.1)
%        Zeile 903  \label{sec:eda}              nach
%                   \subsection{Explorative Analyse und Datenqualitaet} (4.2)
%      main.tex fuehrt jetzt 13 sec-Labels. Alle \ref dieses Abschnitts
%      loesen auf: sec:modellbewertung (2.6), sec:anwendungsfall (3.1),
%      sec:eda (4.2), sec:integration (5.2), sec:rahmen-baselines (5.4).
%      HINWEIS: 2.6 ist noch nicht geschrieben. Der Verweis kompiliert, aber
%      der Grundsatz muss dort auch tatsaechlich stehen.
%   3) Prozentzeichen ist als "\,\%" gesetzt. Falls Kapitel 4 anders setzt:
%      einmal suchen und ersetzen.
%
% DIE DREI ZAHLEN IN ABSATZ 2 SIND NACHGEZAEHLT
%   63  Felder fuehrt das Feldverzeichnis der Quelle
%       (docs/FIR-0001_DataDictionary_fire-incidents.xlsx)
%   15  Felder fordert der Abzug an (prep/s1_daten.py, QUELLEN["SFFD"])
%   10  davon sind gesperrt (prep/config.py, ERGEBNISVARIABLEN)
%   Die Angabe "23 Einsatzspalten, davon bleiben 6" aus der Schreibanleitung
%   ist nicht reproduzierbar und wird nicht verwendet.
% =========================================================================


% ---------------------------- AB HIER KOPIEREN ---------------------------

\subsection{Datenauswahl und Bereinigung}
\label{sec:auswahl-bereinigung}

Der Rohbestand aus Abschnitt~\ref{sec:datengrundlage} umfasst 720.258
Einsatzmeldungen aus 41 Analysis Neighborhoods. Nicht alles davon darf und
nicht alles kann in die Analyse eingehen. Zwei Arten von Entscheidungen führen
von diesem Bestand zum Rahmen der Analyse. Die erste betrifft die inhaltliche
Zulässigkeit einer Größe und folgt dem Grundsatz aus
Abschnitt~\ref{sec:modellbewertung}, dass ein Prädiktor zum Prognosezeitpunkt
tatsächlich verfügbar gewesen sein muss. Die zweite grenzt den
Untersuchungsgegenstand zeitlich und räumlich ab und stützt sich auf die
Qualitätsbefunde aus Abschnitt~\ref{sec:eda}. Beide sind nicht über den
Verarbeitungscode verteilt, sondern als benannte Konstanten in
\texttt{prep/config.py} hinterlegt; jede Auswahlentscheidung ist damit an einem
Ort nachlesbar und an einem Ort änderbar. Am Ende steht der Rahmen, in dem
gerechnet wird: 35 Stadtteile über 132 Monate.

In die Analysedatensätze gelangen ausschließlich Größen, die zum Zeitpunkt des
Alarms bereits feststanden. Kaufman, Rosset und Perlich bezeichnen diese
Anforderung als no-time-machine requirement: Ein zulässiges Modell darf allein
auf Information aufbauen, die zeitlich vor der Zielgröße
liegt.\footcite[559]{Kaufman2011} Die Einsatztabelle führt daneben eine Reihe
von Größen, die den Verlauf und das Ergebnis eines Einsatzes beschreiben, etwa
den geschätzten Sachschaden, die eingesetzten Kräfte oder die Antwortzeit. Sie
hängen eng mit der Einsatzlast zusammen, und gerade darin liegt ihre Gefahr:
Ein Modell, das sie verwendet, erreichte hohe Gütewerte, ohne eine Prognose zu
leisten. Der Ausschluss erfolgt deshalb nicht erst bei der Modellierung,
sondern bereits beim Aufbau des Datensatzes; dieselben Autoren bezeichnen diese
Trennung von Datenhaltung und Vorhersage als learn-predict
separation.\footcite[560]{Kaufman2011} Was nicht in die Datei gelangt, kann
später nicht versehentlich in ein Modell geraten. Von den 63 Feldern des
Feldverzeichnisses fordert der Abzug 15 an, von denen zehn als Ergebnisgrößen
gesperrt sind; dass keine von ihnen den Weg in die Analysedatensätze findet,
prüft \texttt{test\_keine\_ergebnisvariablen} an den fertigen Dateien.

Der verbleibende Bestand wird an drei Stellen bereinigt. Die 269 doppelten
Einsatznummern, die Abschnitt~\ref{sec:eda} ausweist, werden entfernt. Die
Antwortzeit wird auf ein Fenster von null bis
sechzig Minuten geprüft, wobei diese Grenzen eine gesetzte
Plausibilitätsannahme sind und nicht aus den Daten folgen; auf dem verwendeten
Datenstand liegt keine Beobachtung außerhalb des Fensters, sodass die Prüfung
den Bestand nicht verändert. Baujahre außerhalb des Bereichs von 1800 bis 2025
werden als fehlend markiert und nicht gelöscht, weil die übrigen Angaben
derselben Parzelle brauchbar bleiben. An der Zielgröße selbst wird nichts
gekappt, gewinsorisiert oder getrimmt: Die Bereinigung betrifft ausschließlich
Platzhalterwerte der Quellen.

Sechs Stadtteile scheiden ganz aus. Sie sind in Abschnitt~\ref{sec:eda}
benannt; drei sind Parkgebiete ohne nennenswerte Wohnbevölkerung, in denen jede
Pro-Kopf-Größe beliebig groß wird, drei weisen keine durchgängige Abdeckung
durch die ACS-Jahrgänge auf. Entscheidend ist die Form dieses Ausschlusses:
Entfernt werden vollständige Analyseeinheiten, nicht einzelne Zeilen. Ein
zeilenweises Verwerfen unvollständiger Beobachtungen erzeugte ein
unbalanciertes Panel, in dem ein Stadtteil mitten in der Zeitreihe hinzuträte;
Summen und Mittelwerte im Testfenster verschöben sich dann allein durch diesen
Zutritt. Ein rechteckiges Panel ist Voraussetzung dafür, dass die in
Abschnitt~\ref{sec:rahmen-baselines} gebildeten Folds miteinander vergleichbar
sind.

Der Analysezeitraum umfasst die 132 Monate von Januar 2015 bis Dezember 2025.
Sein Beginn ist keine Setzung, sondern eine Folge der Publikationsrealität des
ACS. Einem Einsatz wird nur ein Jahrgang zugeordnet, der spätestens im Vorjahr
erschienen ist (Abschnitt~\ref{sec:integration}); bei den verfügbaren Jahrgängen
2009, 2014, 2019, 2021 und 2023 erhielte ein Einsatz des Jahres 2014 damit den
Jahrgang 2009, und dieser führt die Tabelle B15003 nicht, aus der die
Akademikerquote stammt (Abschnitt~\ref{sec:eda}). Erst ab dem Jahr 2015 greift
der Jahrgang 2014. Das Ende des Zeitraums ist das letzte vollständige
Kalenderjahr. Beide Grenzen sind feste Konstanten und werden nicht aus dem
Datenbestand abgeleitet; andernfalls verschöbe sich der Zeitraum mit jedem
neuen Abzug, und ein angebrochener Randmonat geriete unbemerkt in das
Testfenster. Eine Warnung meldet Monate, deren Aufkommen weniger als die Hälfte
des Medians erreicht; sie filtert nicht, sondern macht auf den Fall aufmerksam.

% ---------------------------- BIS HIER KOPIEREN --------------------------
